\documentclass[10pt]{article}

\usepackage[margin=0.72in]{geometry}
\usepackage{amsmath}
\usepackage{booktabs}
\usepackage[table]{xcolor}
\usepackage{float}
\usepackage{needspace}
\usepackage{microtype}
\usepackage{times}
\usepackage[hidelinks]{hyperref}

\definecolor{headerbg}{HTML}{E9F0EC}
\newcommand{\pp}[1]{\,({#1})}

\title{CRAFT: Supplementary Experimental Analyses}
\author{}
\date{}

\begin{document}
\maketitle

\noindent
This supplement reports analyses that are useful for assessing robustness and
systems cost but are not required for the main argument.  Appendix~A tests
whether duplicate instances within the evaluation benchmark change the
conclusions of the Qwen3-8B training experiments.  Appendix~B characterizes the computational profile of the
evaluation used for Qwen3-8B and all of its training variants.  All
verification rates are percentages.

\appendix
\section{Robustness to Benchmark Deduplication}
\label{app:supp-dedup}

\subsection{Question and protocol}

The full benchmark contains 832 programs.  A separate audit identifies 124
evaluation instances as duplicates of other instances within this benchmark;
removing those redundant instances leaves 708 programs.  This is an
evaluation-set deduplication analysis, distinct from the training--evaluation
overlap audit reported in the main paper.  We use the reduced set as a
sensitivity analysis rather than as a replacement benchmark.  The model
checkpoints, saved response order, sampling budgets
\(k\in\{1,4,8,16,32\}\), composition procedure, and verifier are unchanged.
Thus, differences below reflect the changed evaluation population, not
retraining or resampling.

We report two complementary comparisons.  The \emph{training-stage
comparison} follows representative Bare, RL-from-Bare, and SFT+RL
checkpoints.  The \emph{matched reward ablation} separately compares Binary,
Whole-rollout, Clause-decomposed, and Full reward runs within a common
initialization.  These are distinct evaluation groups and should not be
cross-compared as if they were the same checkpoint.

\subsection{Training-stage comparison}

Tables~\ref{tab:supp-stage-pass} and~\ref{tab:supp-stage-compose} pair the full
and internally deduplicated evaluations.  Parenthesized entries on the 708-program rows
are percentage-point changes from the corresponding 832-program result.

\begin{table}[H]
\centering
\small
\caption{Qwen3-8B training-stage comparison under pass@\(k\).}
\label{tab:supp-stage-pass}
\begin{tabular}{lrrrrrr}
\toprule
\rowcolor{headerbg}
Setting & \(N\) & @1 & @4 & @8 & @16 & @32 \\
\midrule
Bare         & 832 & 8.27 & 16.45 & 20.53 & 24.36 & 28.12 \\
             & 708 & 7.97\pp{-0.30} & 15.87\pp{-0.58} & 19.73\pp{-0.81} & 23.29\pp{-1.07} & 26.84\pp{-1.29} \\
\addlinespace
RL from Bare & 832 & 2.88 & 5.41 & 7.81 & 9.62 & 11.06 \\
             & 708 & 2.54\pp{-0.34} & 4.24\pp{-1.17} & 6.64\pp{-1.17} & 7.91\pp{-1.71} & 9.18\pp{-1.88} \\
\addlinespace
SFT+RL       & 832 & 32.93 & 46.88 & 51.44 & 53.37 & 55.89 \\
             & 708 & 30.37\pp{-2.57} & 43.64\pp{-3.23} & 47.88\pp{-3.56} & 49.72\pp{-3.65} & 52.40\pp{-3.49} \\
\bottomrule
\end{tabular}
\end{table}

\begin{table}[H]
\centering
\small
\caption{Qwen3-8B training-stage comparison under compose@\(k\).}
\label{tab:supp-stage-compose}
\begin{tabular}{lrrrrrr}
\toprule
\rowcolor{headerbg}
Setting & \(N\) & @1 & @4 & @8 & @16 & @32 \\
\midrule
Bare         & 832 & 43.75 & 55.17 & 59.01 & 61.18 & 62.50 \\
             & 708 & 42.09\pp{-1.66} & 53.81\pp{-1.35} & 57.91\pp{-1.10} & 60.17\pp{-1.01} & 61.44\pp{-1.06} \\
\addlinespace
RL from Bare & 832 & 51.80 & 60.46 & 64.18 & 68.87 & 70.67 \\
             & 708 & 48.87\pp{-2.93} & 57.20\pp{-3.25} & 61.16\pp{-3.02} & 65.68\pp{-3.19} & 67.51\pp{-3.16} \\
\addlinespace
SFT+RL       & 832 & 66.23 & 72.48 & 74.40 & 75.12 & 75.60 \\
             & 708 & 62.99\pp{-3.23} & 69.07\pp{-3.41} & 71.05\pp{-3.35} & 71.89\pp{-3.23} & 72.32\pp{-3.28} \\
\bottomrule
\end{tabular}
\end{table}

\Needspace{8\baselineskip}
\paragraph{Stage conclusions.}
Removing within-benchmark duplicates lowers the reported rates but preserves the
shape of both curves.  Across the two tables, the absolute changes range from
0.30 to 3.65 points.  The larger shifts for SFT+RL indicate that the removed
subset is comparatively favorable to that checkpoint, but SFT+RL remains the
strongest stage at every budget under both pass and composition.  Likewise,
RL from Bare continues to improve compose@\(k\) over Bare while reducing
pass@\(k\), preserving the distinction between standalone-response quality
and cross-response clause utility.

\subsection{Stability of the reward ablation}

The central robustness question is whether internal benchmark deduplication changes the reward
ranking.  Tables~\ref{tab:supp-reward-compose} and
\ref{tab:supp-reward-pass} report the complete matched ablation.  Each reward
has adjacent 832- and 708-program rows.  ``Bare (reference)'' appears only in
the Bare-initialization panel; it is not an additional trained reward.

\begin{table}[H]
\centering
\small
\caption{Reward ablation under compose@\(k\).  The two panels distinguish RL
directly from Bare from RL after SFT.}
\label{tab:supp-reward-compose}
\begin{tabular}{lrrrrrr}
\toprule
\rowcolor{headerbg}
Reward & \(N\) & @1 & @4 & @8 & @16 & @32 \\
\midrule
\multicolumn{7}{l}{\textit{RL directly from Bare}} \\
Bare (reference)  & 832 & 43.75 & 55.17 & 59.01 & 61.18 & 62.50 \\
                  & 708 & 42.09 & 53.81 & 57.91 & 60.17 & 61.44 \\
\addlinespace
Binary            & 832 & 0.00 & 0.00 & 0.00 & 0.00 & 0.00 \\
                  & 708 & 0.00 & 0.00 & 0.00 & 0.00 & 0.00 \\
\addlinespace
Whole-rollout     & 832 & 20.79 & 21.88 & 22.24 & 22.84 & 23.32 \\
                  & 708 & 19.63 & 20.76 & 21.19 & 21.61 & 22.03 \\
\addlinespace
Clause-decomposed & 832 & 48.80 & 57.69 & 60.70 & 63.70 & 66.35 \\
                  & 708 & 46.75 & 55.93 & 59.04 & 61.86 & 64.69 \\
\addlinespace
Full (ours)       & 832 & 51.80 & 60.46 & 64.18 & 68.87 & 70.67 \\
                  & 708 & 48.87 & 57.20 & 61.16 & 65.68 & 67.51 \\
\midrule
\multicolumn{7}{l}{\textit{RL after SFT}} \\
Binary            & 832 & 7.09 & 10.82 & 13.34 & 14.54 & 15.99 \\
                  & 708 & 6.78 & 10.03 & 12.57 & 13.84 & 15.54 \\
\addlinespace
Whole-rollout     & 832 & 48.80 & 56.37 & 60.22 & 63.10 & 66.59 \\
                  & 708 & 47.03 & 53.67 & 57.06 & 59.89 & 63.42 \\
\addlinespace
Clause-decomposed & 832 & 67.91 & 75.84 & 76.80 & 77.64 & 78.25 \\
                  & 708 & 64.69 & 72.88 & 73.73 & 74.72 & 75.14 \\
\addlinespace
Full (ours)       & 832 & 68.51 & 75.72 & 77.52 & 78.25 & 78.37 \\
                  & 708 & 64.83 & 72.46 & 74.58 & 75.28 & 75.42 \\
\bottomrule
\end{tabular}
\end{table}

\begin{table}[H]
\centering
\small
\caption{Reward ablation under pass@\(k\).  Pass measures whether any complete
response verifies alone and therefore need not follow the compose ranking.}
\label{tab:supp-reward-pass}
\begin{tabular}{lrrrrrr}
\toprule
\rowcolor{headerbg}
Reward & \(N\) & @1 & @4 & @8 & @16 & @32 \\
\midrule
\multicolumn{7}{l}{\textit{RL directly from Bare}} \\
Bare (reference)  & 832 & 8.27 & 16.45 & 20.53 & 24.36 & 28.12 \\
                  & 708 & 7.97 & 15.87 & 19.73 & 23.29 & 26.84 \\
\addlinespace
Binary            & 832 & 4.45 & 4.45 & 4.45 & 4.45 & 4.45 \\
                  & 708 & 4.52 & 4.52 & 4.52 & 4.52 & 4.52 \\
\addlinespace
Whole-rollout     & 832 & 15.50 & 16.59 & 16.83 & 17.67 & 18.39 \\
                  & 708 & 15.54 & 16.67 & 16.95 & 17.66 & 18.08 \\
\addlinespace
Clause-decomposed & 832 & 9.01 & 15.87 & 18.87 & 20.31 & 21.63 \\
                  & 708 & 8.19 & 14.69 & 17.66 & 19.21 & 20.62 \\
\addlinespace
Full (ours)       & 832 & 2.88 & 5.41 & 7.81 & 9.62 & 11.06 \\
                  & 708 & 2.54 & 4.24 & 6.64 & 7.91 & 9.18 \\
\midrule
\multicolumn{7}{l}{\textit{RL after SFT}} \\
Binary            & 832 & 8.77 & 11.30 & 13.22 & 14.30 & 15.50 \\
                  & 708 & 8.47 & 10.73 & 12.57 & 13.70 & 15.11 \\
\addlinespace
Whole-rollout     & 832 & 41.11 & 49.28 & 53.12 & 55.53 & 58.29 \\
                  & 708 & 38.56 & 46.05 & 49.58 & 51.84 & 54.66 \\
\addlinespace
Clause-decomposed & 832 & 35.94 & 51.44 & 54.33 & 57.69 & 61.30 \\
                  & 708 & 32.06 & 47.32 & 50.56 & 53.39 & 56.78 \\
\addlinespace
Full (ours)       & 832 & 36.78 & 51.92 & 56.49 & 59.50 & 61.54 \\
                  & 708 & 32.91 & 47.88 & 52.40 & 55.23 & 57.34 \\
\bottomrule
\end{tabular}
\end{table}

\paragraph{The composition result is robust.}
Clause-level rewards dominate response-level rewards under compose@\(k\) on
both evaluation sets and from both initializations.  At compose@32 the strict
ordering is unchanged:
\[
\text{Full} > \text{Clause-decomposed} > \text{Whole-rollout} > \text{Binary}.
\]
After SFT, Full and Clause-decomposed are close: Clause-decomposed is slightly
higher at @4, whereas Full is higher at @1, @8, @16, and @32 on both sets.
These small local crossings do not alter the endpoint ranking or the larger
gap between clause-level and response-level credit.

\paragraph{Pass and compose expose different effects.}
Whole-rollout can retain stronger pass@\(k\) than the clause-level objectives,
especially from Bare, while performing substantially worse under composition.
This is consistent with the objectives: response-level credit favors complete
standalone answers, whereas clause-level credit preserves useful partial
content for later composition.  The within-benchmark deduplication control
retains this contrast.

\subsection{Negative-trace sampling ablation}

We next isolate negative-trace sampling while holding the Full reward fixed.
Structure uses the target-hidden negative-trace construction defined in the
main paper; Random replaces it with a budget-matched random construction.  We
compare both samplers under RL and SFT+RL.  As in the main paper, RL denotes
training directly from Bare, whereas SFT+RL starts from the SFT model.

\begin{table}[htbp]
\centering
\caption{Qwen3-8B negative-trace sampling ablation on all 832 programs
under pass@\(k\), with the Full reward.  Difference rows report Structure
minus Random in percentage points.}
\label{tab:sampler-pass-all}
\small
\begin{tabular}{llrrrrr}
\toprule
Training & Sampler & @1 & @4 & @8 & @16 & @32 \\
\midrule
SFT+RL & Structure & 37.98 & 51.80 & 56.13 & 59.98 & 62.74 \\
       & Random    & 37.26 & 48.20 & 53.85 & 58.05 & 60.70 \\
       & $\Delta$ & +0.72 & +3.61 & +2.28 & +1.92 & +2.04 \\
\addlinespace
RL     & Structure & 8.17 & 13.10 & 14.90 & 17.43 & 19.35 \\
       & Random    & 5.65 & 11.18 & 13.22 & 15.87 & 18.03 \\
       & $\Delta$ & +2.52 & +1.92 & +1.68 & +1.56 & +1.32 \\
\bottomrule
\end{tabular}
\end{table}

\begin{table}[htbp]
\centering
\caption{Qwen3-8B negative-trace sampling ablation on all 832 programs
under compose@\(k\), with the Full reward.  Difference rows report Structure
minus Random in percentage points.}
\label{tab:sampler-compose-all}
\small
\begin{tabular}{llrrrrr}
\toprule
Training & Sampler & @1 & @4 & @8 & @16 & @32 \\
\midrule
SFT+RL & Structure & 67.55 & 76.08 & 77.52 & 78.61 & 78.73 \\
       & Random    & 68.39 & 75.24 & 77.16 & 78.25 & 78.49 \\
       & $\Delta$ & -0.84 & +0.84 & +0.36 & +0.36 & +0.24 \\
\addlinespace
RL     & Structure & 49.28 & 57.93 & 62.26 & 64.90 & 66.83 \\
       & Random    & 50.48 & 59.86 & 62.50 & 64.78 & 66.11 \\
       & $\Delta$ & -1.20 & -1.93 & -0.24 & +0.12 & +0.72 \\
\bottomrule
\end{tabular}
\end{table}

Tables~\ref{tab:sampler-pass-strata} and
\ref{tab:sampler-compose-strata} disaggregate the same comparison by program
class.  The ``All'' results are reported separately above to avoid repeating
the aggregate in the stratified tables.

\begin{table}[htbp]
\centering
\caption{Pass@\(k\) negative-trace sampling ablation by program class.}
\label{tab:sampler-pass-strata}
\scriptsize
\setlength{\tabcolsep}{5pt}
\begin{tabular}{llrrrrr}
\toprule
Training & Sampler & @1 & @4 & @8 & @16 & @32 \\
\midrule
\multicolumn{7}{l}{\textit{Linear ($n=316$)}} \\
SFT+RL & Structure & 40.51 & 56.33 & 59.18 & 64.56 & 67.09 \\
       & Random    & 37.66 & 49.68 & 56.01 & 59.49 & 63.61 \\
RL     & Structure & 10.13 & 13.29 & 15.19 & 18.35 & 20.25 \\
       & Random    & 8.23 & 13.29 & 15.19 & 17.41 & 19.62 \\
\addlinespace
\multicolumn{7}{l}{\textit{NLA ($n=50$)}} \\
SFT+RL & Structure & 6.00 & 6.00 & 8.00 & 8.00 & 8.00 \\
       & Random    & 6.00 & 8.00 & 8.00 & 8.00 & 8.00 \\
RL     & Structure & 0.00 & 0.00 & 2.00 & 2.00 & 2.00 \\
       & Random    & 0.00 & 0.00 & 2.00 & 2.00 & 2.00 \\
\addlinespace
\multicolumn{7}{l}{\textit{Loopy ($n=466$)}} \\
SFT+RL & Structure & 39.70 & 53.65 & 59.23 & 62.45 & 65.67 \\
       & Random    & 40.34 & 51.50 & 57.30 & 62.45 & 64.38 \\
RL     & Structure & 7.73 & 14.38 & 16.09 & 18.45 & 20.60 \\
       & Random    & 4.51 & 10.94 & 13.09 & 16.31 & 18.67 \\
\bottomrule
\end{tabular}
\end{table}

\begin{table}[htbp]
\centering
\caption{Compose@\(k\) negative-trace sampling ablation by program class.}
\label{tab:sampler-compose-strata}
\scriptsize
\setlength{\tabcolsep}{5pt}
\begin{tabular}{llrrrrr}
\toprule
Training & Sampler & @1 & @4 & @8 & @16 & @32 \\
\midrule
\multicolumn{7}{l}{\textit{Linear ($n=316$)}} \\
SFT+RL & Structure & 69.94 & 79.11 & 80.70 & 81.33 & 81.33 \\
       & Random    & 72.47 & 78.48 & 80.70 & 81.33 & 81.33 \\
RL     & Structure & 53.16 & 62.97 & 66.77 & 70.25 & 71.20 \\
       & Random    & 52.53 & 62.34 & 65.19 & 66.77 & 68.99 \\
\addlinespace
\multicolumn{7}{l}{\textit{NLA ($n=50$)}} \\
SFT+RL & Structure & 12.00 & 14.00 & 14.00 & 14.00 & 14.00 \\
       & Random    & 12.00 & 12.00 & 12.00 & 12.00 & 12.00 \\
RL     & Structure & 10.00 & 12.00 & 12.00 & 12.00 & 12.00 \\
       & Random    & 12.00 & 12.00 & 12.00 & 14.00 & 14.00 \\
\addlinespace
\multicolumn{7}{l}{\textit{Loopy ($n=466$)}} \\
SFT+RL & Structure & 71.89 & 80.69 & 82.19 & 83.69 & 83.91 \\
       & Random    & 71.67 & 79.83 & 81.76 & 83.26 & 83.69 \\
RL     & Structure & 50.86 & 59.44 & 64.59 & 66.95 & 69.74 \\
       & Random    & 53.22 & 63.30 & 66.09 & 68.88 & 69.74 \\
\bottomrule
\end{tabular}
\end{table}

\paragraph{Negative-trace sampling mainly affects standalone success.}
On the full benchmark, Structure improves pass@\(k\) over Random at every
response budget: the gain ranges from 0.72 to 3.61 points under SFT+RL and
from 1.32 to 2.52 points under RL.  The stratified results
locate most of this gain in Linear and Loopy
programs.  The small NLA subset changes by at
most two points, corresponding to a single program, and often ties.

\paragraph{Composition is insensitive to negative-trace sampling.}
For compose@\(k\), the aggregate Structure--Random difference lies between
-1.93 and +0.84 points and changes sign with the training stage and response
budget.  The class-level results show the same mixed pattern: neither sampler
dominates uniformly across Linear, NLA, and Loopy programs.  Thus, Structure
provides a modest and consistent benefit for standalone
pass@\(k\), but has no systematic effect on compose@\(k\).  This ablation
also separates negative-trace sampling from the
larger reward-design effects above: the qualitative advantage of clause-level
credit under composition cannot be attributed to the sampling strategy.

\section{Cost and Systems Profile of Qwen3-8B and Its Training Variants}
\label{app:supp-resources}

\subsection{Measurement scope}

We profile the common evaluation configuration used for the Qwen3-8B base
model and all of its training variants.  Each checkpoint is evaluated on all
832 programs with eight rollouts per task.  The reported totals are therefore
the cost of one complete checkpoint evaluation, not the sum across all
variants.  The measurement boundary includes prompt processing, response
generation, candidate extraction and pooling, Houdini filtering, and final
Frama-C/WP verification.  Token totals, verifier timings, and memory
measurements are recorded quantities.  Multi-GPU generation time and the
80-worker wall-clock decomposition use the measured throughput and aggregate
work, and therefore describe this serving configuration rather than a
hardware-independent latency guarantee.

\subsection{Generation workload}

\begin{table}[H]
\centering
\small
\caption{Token consumption for one Qwen3-8B checkpoint or training-variant
evaluation over 832 programs.  Each task uses eight rollouts; per-task values
are rounded to the nearest token.}
\label{tab:supp-resource-tokens}
\begin{tabular}{lrr}
\toprule
\rowcolor{headerbg}
Component & Mean per task & Total over 832 tasks \\
\midrule
Prompt tokens              & 1,113  & 925,635 \\
Output tokens (8 rollouts) & 11,057 & 9,199,244 \\
Total tokens               & 12,169 & 10,124,879 \\
\bottomrule
\end{tabular}
\end{table}

Each rollout produces 1,382 output tokens on average and yields approximately
19.6 parsed clauses.  Pooling and deduplicating the first eight rollouts leaves
58.8 candidate invariants per task on average, with a maximum of 106.  Thus,
candidate deduplication removes substantial textual repetition before formal
verification: eight responses yield roughly 157 parsed clauses in total, but
only 58.8 distinct candidates enter the pooled set on average.

\begin{table}[H]
\centering
\small
\caption{Measured vLLM output throughput and derived generation time.}
\label{tab:supp-resource-generation}
\begin{tabular}{lr}
\toprule
\rowcolor{headerbg}
Configuration & Throughput or time \\
\midrule
A800-80G throughput & 5,833 output tokens/s/GPU \\
One-GPU generation & 26.3 min \\
Four-GPU generation & 6.6 min \\
\bottomrule
\end{tabular}
\end{table}

The one-GPU time follows from the 9,199,244 generated tokens and measured
output-token throughput.  The four-GPU value assumes even task allocation at
the same per-GPU throughput.  Generation is therefore inexpensive relative to
the subsequent proof workload once multiple GPUs are available.

\subsection{Formal verification dominates compute}

\begin{table}[H]
\centering
\small
\caption{Frama-C/Houdini verification cost overall and by benchmark subset.}
\label{tab:supp-resource-verification}
\begin{tabular}{lrrrrrr}
\toprule
\rowcolor{headerbg}
Subset & \(N\) & Mean & Median & P95 & CPU time & Budget hits \\
\midrule
All    & 832 & 123.4 s & 25.5 s  & 551.1 s & 28.5 h & 55 (6.6\%) \\
Linear & 316 & 112.7 s & 24.6 s  & 530.7 s & 9.9 h  & -- \\
NLA    & 50  & 292.0 s & 252.6 s & 866.8 s & 4.1 h  & -- \\
Loopy  & 466 & 112.6 s & 22.3 s  & 564.5 s & 14.6 h & -- \\
\bottomrule
\end{tabular}
\end{table}

Verification time is strongly right-skewed: the overall median is 25.5
seconds, while the P95 is 551.1 seconds and 55 tasks hit the 300-second proof
budget.  NLA is the most expensive subset per task.  Although it contains only
50 programs, it consumes 4.1 CPU-hours and averages 292.0 seconds (4.9
minutes) per program.  Loopy consumes the most aggregate CPU time because it
is the largest subset, not because its per-task mean exceeds Linear.

\subsection{Memory and parallel execution}

\begin{table}[H]
\centering
\small
\caption{Memory footprint and verifier concurrency.}
\label{tab:supp-resource-memory}
\begin{tabular}{ll}
\toprule
\rowcolor{headerbg}
Component & Memory or concurrency \\
\midrule
GPU model weights          & 15.27 GiB \\
GPU KV cache               & 55.43 GiB (capacity: approximately 403K tokens) \\
Total per GPU              & approximately 73.6/80 GiB \\
CPU per Frama-C worker     & approximately 200--500 MB \\
Concurrent CPU workers     & 60--100 \\
\bottomrule
\end{tabular}
\end{table}

The 80-GiB GPU is close to capacity, but most allocated memory is the KV cache
rather than model weights.  On the verifier side, 80 concurrent workers imply
approximately 16--40 GB of worker memory from the observed per-process range,
excluding the operating system and orchestration overhead.  The pipeline is
therefore GPU-memory-bound during serving and CPU-parallelism-bound during
formal verification.

\subsection{End-to-end bottleneck}

\begin{table}[H]
\centering
\small
\caption{End-to-end wall-clock decomposition with four GPUs and 80 CPU
verification workers.}
\label{tab:supp-resource-wallclock}
\begin{tabular}{lrr}
\toprule
\rowcolor{headerbg}
Stage & Wall-clock time & Share \\
\midrule
LLM generation       & 6.6 min  & 23\% \\
Frama-C verification & 21.4 min & 77\% \\
Total                & approximately 28 min & -- \\
\bottomrule
\end{tabular}
\end{table}

The systems result mirrors the CPU accounting: proof is the bottleneck.
Before parallelization, generation contributes only 1.5\% of the summed stage
times (1,577 seconds of generation versus 102,675 CPU-seconds of
verification).  Parallel execution narrows this imbalance, but verification
still occupies 77\% of the estimated end-to-end wall-clock time.  Additional
generation throughput alone would therefore have limited impact; reducing
proof obligations, avoiding repeated verifier work, or improving verifier
load balancing targets the dominant cost.

\end{document}
